\chapter{考察}
\label{chap:考察}

\section{Phase 1 仮説（言語-動作整合潜在空間の獲得）の検証}
\label{sec:contrastive_discussion}

Phase 1 では，オノマトペのテキスト埋め込みと歩容モーション埋め込みを
同一の潜在空間上で対応付け，検索タスクを通じて言語-動作の整合を獲得できるかを検証した．

第\ref{chap:結果}章の結果から，最も性能に影響したのは温度パラメータ $\tau$ であった．
低温側（$\tau=0.02$〜$0.05$）で検索性能が安定する一方，
$\tau=0.10$ まで上げると性能が明確に落ちた．
これは，温度が高いと類似度分布が平坦化し，
正例と負例の区別がつきにくくなるためと考えられる．
また，VAE の正則化項を適度に加えることで汎化性能が改善する傾向が見られた一方，
コントラスト損失の重みを強くしすぎると学習が不安定になり，
かえって性能が低下した．

一方で，潜在空間の構造と検索性能の間にはギャップがある．
図\ref{fig:contrastive_pca_fine}の PCA 可視化では
オノマトペ間でクラスタの重なりが残り，シルエット係数も負の値を示しており，
幾何学的にクラスタが明確に分離されているとは言えない．
それにもかかわらず，テキストからモーションを検索する t2m 指標では一定の性能が得られている．
これは，クラスタが完全に分離していなくても，
テキスト側の埋め込みが比較的安定しているために，
そこを基準とした検索で語彙と歩容の対応が成立していると考えられる．
一方で，モーション側からテキストを検索する m2t の性能が t2m より低い点は，
モーション埋め込みの分散が大きく，逆方向の検索が難しいことを反映している．

以上より，Phase 1 の仮説は言語-動作の対応付けという意味では支持されるが，
潜在空間の明確な分離や双方向検索の対称性までは達成できておらず，課題が残る．

\section{Phase 1 の達成が限定的であった要因}
\label{sec:phase1_factor_discussion}

Phase 1 で残された課題の背景には，主に以下の 4 つの要因がある．

\subsection{データ規模と温度感度}
本研究の訓練データは約 222 サンプルと小規模であり，
コントラスト学習が本来想定するデータ規模から大きく離れている．
小規模データでは，温度を低くしてコントラストを強くとる方が
各クラスの代表的な特徴を捉えやすい反面，
温度設定のわずかな違いが性能に直結してしまう．
実際に，$\tau$ を少し変えただけで結果が不安定になる傾向が確認された．

\subsection{損失関数間のバランス}
VAE の再構成損失は潜在空間の滑らかさを保ち，汎化を助ける役割を果たす．
しかし，コントラスト損失を強くしすぎると，
各クラスを無理に引き離そうとする力が強まり，
学習が不安定化する．
つまり，「意味的に整合した空間を作る」ことと
「潜在表現を滑らかに保つ」ことの間に本質的なトレードオフがあり，
そのバランス調整が Phase 1 のボトルネックとなっている．

\subsection{クラスタ未分離でも検索が成立する理由}
PCA 上ではクラスタが重なっているにもかかわらず検索が機能しているのは，
本手法がクラスタを完全に分離することではなく，
各クラスの代表点に対して相対的にどれが近いかで順位を付けているためだと考えられる．
ただし，この仕組みではクラス境界付近のサンプルで取り違えが起きやすく，
実際に意味の近い語彙同士（例：「すたすた」と「てくてく」）では
混同が生じている．

\subsection{テキストとモーションの表現の非対称性}
テキスト埋め込みは事前学習済み言語モデルに由来するため比較的安定しており，
検索のアンカーとして機能しやすい．
一方で，モーション埋め込みは歩容データから新たに学習するため分散が大きく，
とくにモーション側を起点とする検索（m2t）で性能が落ちやすい．
この非対称性は，双方向で同等の対応付けを実現するうえでの構造的な課題である．

以上の要因が相互に作用した結果として，Phase 1 は
限られたデータ・限られたハイパーパラメータ探索の中で
一定の整合を獲得しつつも，完全な分離・対称性には至らなかった．

\section{Phase 2 仮説（オノマトペ条件付き歩行生成）の検証}
\label{sec:rl_discussion}

Phase 2 では，Phase 1 で獲得した潜在空間を利用し，
オノマトペを条件として与えることで，ロボットの歩行スタイルを語彙に応じて切り替えられるかを検証した．

第\ref{subsec:rl_results}節の結果から，語彙間で速度やコサイン類似度に差が出ており，
条件付けが歩行に影響を与えていること自体は確認できた．
なお，本研究では速度条件を足踏み回避の補助として導入しているため，
主な評価基準は速度の一致ではなく，語彙ごとの相対的な速度差や見た目の妥当性である．
ただし，その基準でも姿勢の品質や語彙間の分離は十分とは言えなかった．

条件付けの経路自体が機能しているかを確かめるため，
「すたすた」と「のろのろ」の 2 語のみに絞った補助実験を行った
（表\ref{tab:rl_two_style_ablation_stylew4_cmd03}）．

\begin{table}[H]
  \centering
  \caption{2語補助評価（style\_weight=4.0，速度指令値 0.3）}
  \label{tab:rl_two_style_ablation_stylew4_cmd03}
  \begin{tabular}{lccccc} \toprule
    \textbf{オノマトペ} & $\mathbf{v_x}$ & \textbf{CosC} & \textbf{style score} & \textbf{Err} & \textbf{Fall} \\ \midrule
    すたすた & 0.3436 & 0.3277 & -0.1051 & 4.3016 & 0.0000 \\
    のろのろ & 0.2843 & 0.6077 & +0.1143 & 3.4666 & 0.0000 \\ \bottomrule
  \end{tabular}
\end{table}

この結果では，コサイン類似度とスタイルスコアに 2 語間で差が出ており，
条件付けの入力がまったく無視されているわけではない．
また，この 2 語設定での類似度や関節誤差は，
11 語を同時に条件付けた主比較実験（C 条件，表\ref{tab:rl_cd_onom_dynamics}，表\ref{tab:rl_abcd_onom_similarity}）と
大きく変わらない水準にある．
語彙数を減らしても結果が改善しないことから，
条件付けの経路自体に問題があるというよりも，
報酬の設計や方策の最適化のほうが
歩容差を出しにくくしている要因だと考えられる．

主比較実験（表\ref{tab:rl_abcd_main_m2999}）でも，
スタイル重みを大きくした条件ほど語彙間の速度差が拡大する傾向が確認できる．
ただし，現状ではその差が主に速度の大小として現れている段階であり，
関節運動パターンの質的な違いとして歩容スタイルが分化するには至っていない．
また，C 条件で未知語 7 語を入力した追加評価
（表\ref{tab:rl_c_unknown_onom_velx}）では，
Vel X が 0.4793〜0.5228 に集中しており，
既知語の C 条件で見られた大きな速度差
（0.1721〜0.3179，表\ref{tab:rl_cd_onom_dynamics}）は再現されなかった．
既知語では語彙ごとの速度差が出るのに対し，
未知語ではそうした差が現れにくく，
スタイルの一般化にはまだ至っていないと言える．

\subsection{速度追従結果と HOYO 相対速度差の比較（C 条件）}
\label{subsec:discussion_c_vs_hoyo_speed}

次に，速度追従結果と HOYO 由来の速度差を直接比較する．
語彙間速度差が最も明確に出た C 条件について，
各語彙の速度（表\ref{tab:rl_cd_onom_dynamics}）を
通常の速度で割って相対速度（通常=1）に変換し，
HOYO 側の相対速度（表\ref{tab:hoyo_relative_speed_table}）も
通常を基準に正規化して，両者の絶対差を算出した．

\begin{table}[H]
  \centering
  \caption{C 条件の速度・相対速度と HOYO 相対速度の比較（通常=1）}
  \label{tab:discussion_c_vs_hoyo_relative_speed}
  \small
  \begin{tabular}{lcccc} \toprule
    \textbf{オノマトペ} & \textbf{速度} & \textbf{相対速度} & \textbf{教師の相対速度} & \textbf{差分 $|\Delta|$} \\ \midrule
    通常 & 0.253 & 1.000 & 1.000 & 0.000 \\
    すたすた & 0.318 & 1.257 & 1.121 & 0.136 \\
    せかせか & 0.266 & 1.051 & 1.121 & 0.070 \\
    てくてく & 0.280 & 1.107 & 0.794 & 0.313 \\
    どっしどっし & 0.243 & 0.960 & 0.784 & 0.177 \\
    とぼとぼ & 0.241 & 0.953 & 0.480 & 0.473 \\
    のしのし & 0.172 & 0.680 & 0.573 & 0.107 \\
    のろのろ & 0.273 & 1.079 & 0.502 & 0.577 \\
    ぶらぶら & 0.292 & 1.154 & 0.636 & 0.518 \\
    よたよた & 0.270 & 1.067 & 0.522 & 0.545 \\
    よろよろ & 0.254 & 1.004 & 0.307 & 0.697 \\ \bottomrule
  \end{tabular}
\end{table}

表\ref{tab:discussion_c_vs_hoyo_relative_speed}より，
「すたすた」「のしのし」は HOYO 側と近い一方で，
「よろよろ」「のろのろ」「ぶらぶら」は差が大きい．
速い/遅いの大まかな傾向は部分的に再現されているが，
語彙ごとの細かな相対速度の順序までは再現できていない．

歩容スタイルが十分に分かれなかった要因として，以下の 6 点が考えられる．
\begin{enumerate}
  \item \textbf{報酬による歩容の平均化}：
        現行のスタイル報酬は正例との類似度だけで評価しており，
        どの語彙でもそれなりに高い点を取れる「無難な歩き方」に
        落ち着きやすい構造になっている．
  \item \textbf{教師データ（HOYO）の規模が小さい}：
        本研究で教師として用いる HOYO は語彙数に対してサンプル数が限られており，
        語彙内のばらつきや語彙間境界を十分に学習しにくい．
        その結果，方策側では語彙固有の特徴よりも平均的な歩容へ寄りやすい．
  \item \textbf{2D表現では進行の仕方が直接観測できない}：
        HOYO は正面視点の 2D 関節系列であり，
        「どの程度前に進んだか」を 3D の並進量として直接扱えない．
        そのため，速度や歩幅の差は間接的な見かけ量に依存し，
        スタイル差として安定して学習されにくい．
  \item \textbf{奥行き情報の消失}：
        2D 投影・センタリング処理により奥行き方向の情報が失われるため，
        前後動作の違いや身体の立体的な使い方が潜在表現で弱まりやすい．
        この情報欠落が，語彙ごとの質感差の分離を難しくしている．
  \item \textbf{姿勢の崩れ}：
        一部の試行では上肢の姿勢が崩れており，
        速度に差があっても見た目に「歩き方が違う」とは感じにくかった．
  \item \textbf{分離指標の停滞}：
        マージン混同行列の非対角成分はいずれの条件でもほぼゼロ付近にとどまり
        （表\ref{tab:rl_abcd_main_m2999}），
        学習を進めても語彙間の識別が改善する傾向は見られなかった．
\end{enumerate}

まとめると，条件付け自体は機能しているが，
データ規模と 2D 表現の情報制約に加えて，
最適化の過程で歩容が平均化され，語彙ごとの差が薄れてしまった．
結果として，語彙に応じた速度差や類似度差は出ているものの，
「オノマトペごとに見た目の異なる歩き方になっている」とまでは言えない．

\section{Phase 2 の検証限界と今後の課題}
\label{sec:phase2_limitations}

本研究の Phase 2 の評価は，
条件付けの経路が機能しているかの確認が主であり，
語彙間のスタイル差を定量的に判定する仕組みにはなっていない．

この限界を踏まえると，まずスタイル報酬の設計を見直す必要がある．
現行の報酬は正例との類似度だけを見ているが，
他の語彙（負例）との類似度を下げる項を加えれば，
語彙間の差をより直接的に広げられる可能性がある．
具体的には，従来の報酬
\begin{align}
  r_{\mathrm{legacy}}
  = \beta_{\mathrm{text}}r_{\mathrm{text}}
  + \beta_{\mathrm{teacher}}r_{\mathrm{teacher}}
\end{align}
に対して，現在ラベル以外の最大類似度 \(r_{\mathrm{neg}}\) を導入し，
\begin{align}
  r_{\mathrm{pos}}
  &= \beta_{\mathrm{text}}r_{\mathrm{text}}
  + \beta_{\mathrm{teacher}}r_{\mathrm{teacher}}, \\
  p
  &= \max(0,\;r_{\mathrm{neg}}-r_{\mathrm{teacher}}+m), \\
  r_{\mathrm{hardneg}}
  &= r_{\mathrm{pos}}-\lambda_{\mathrm{neg}}p
\end{align}
とする hard-negative 形式が考えられる．
負例の類似度が教師類似度を上回ったときだけペナルティがかかるため，
歩行の安定性を大きく損なわずに語彙間の分離を促せる見込みがある．
ただし，マージン $m$ やペナルティ重み $\lambda_{\mathrm{neg}}$ の調整は慎重に行う必要がある．

報酬設計に加えて，以下の方向性が今後の改善として有力である．

\begin{enumerate}
  \item \textbf{3D モーションデータへの移行}：
        本研究では HOYO の 2D 正面視点系列を用いたが，
        奥行き情報の消失が歩容差の学習を難しくしていた．
        3D モーションキャプチャデータを教師として使えれば，
        歩幅や体幹の前傾など，2D では観測しにくい特徴を
        潜在空間に直接反映でき，スタイル分離の精度が大きく改善する可能性がある．
  \item \textbf{姿勢品質の直接的な制約}：
        現状では上肢の姿勢崩れがスタイル差の知覚を妨げている．
        教師モーションとの関節角誤差を報酬に組み込み，
        とくに腕や体幹に重みを置くことで，
        速度差だけでなく見た目の質も制御できるようになる．
  \item \textbf{カリキュラム学習}：
        まず安定した前進歩行を獲得してから，
        スタイル報酬の重みを段階的に上げていく構成にすれば，
        転倒リスクを抑えつつ語彙間の差を広げやすくなる．
        現行のランプスケジュールをより精緻に設計する余地がある．
  \item \textbf{主観評価の導入}：
        速度差やコサイン類似度といった定量指標だけでは，
        「歩き方が違って見えるか」を捉えきれない．
        人間による対比較実験を取り入れることで，
        定量指標と知覚的な印象のずれを把握できる．
  \item \textbf{未知語の一般化評価}：
        未知語セットを事前に固定し，既知語と同じ条件で評価する
        プロトコルを整備する必要がある．
        対照学習の潜在空間における未知語の位置と
        実際の歩行出力の対応を分析すれば，
        一般化がどこで途切れているかを特定しやすくなる．
\end{enumerate}

以上のように，Phase 2 の検証は現時点では限定的であり，
仮説をより確かなものにするには評価の枠組み自体の見直しが必要である．
